\section{Periodi}
%%%%%%%%%%%%%%%%%%%%%%%%%%%%%%%%%%%%%%%%%%%%%%%%%%%%%%%%%%%%%%
% Semplicemente un'idea di come potrebbe essere, da valutare %
%%%%%%%%%%%%%%%%%%%%%%%%%%%%%%%%%%%%%%%%%%%%%%%%%%%%%%%%%%%%%%
    \subsection{RTB}
        \subsubsection{I Periodo (dal 2024-10-XY al 2024-10-XY)}
        Punti salienti/Obiettivi: \\
        Rischi incontrati e intervento:
        \subsubsection{II Periodo (dal 2024-10-XY al 2024-11-XY)}
        Punti salienti/Obiettivi: \\
        Rischi incontrati e intervento:
        \subsubsection{III Periodo (dal 2024-11-XY al 2024-11-XY)}
        Punti salienti/Obiettivi: \\
        Rischi incontrati e intervento:
    \subsection{PB}
        \subsubsection{IV Periodo (dal 2024-12-XY al 2025-01-XY)}
        Punti salienti/Obiettivi: \\
        Rischi incontrati e intervento:
        \subsubsection{V Periodo (dal 2025-01-XY al 2025-01-XY)}
        Punti salienti/Obiettivi: \\
        Rischi incontrati e intervento:
        \subsubsection{VI Periodo (dal 2025-01-XY al 2025-02-XY)}
        Punti salienti/Obiettivi: \\
        Rischi incontrati e intervento:
\\[20pts]
%%%%%%%%%%%%%%%%%%%%%%%%%%%%%%%%%%%%%%%%%%%%%%%%%%%%%%%%%
% Qui sotto invece e' riportato quello fatto dai Mattei %
%%%%%%%%%%%%%%%%%%%%%%%%%%%%%%%%%%%%%%%%%%%%%%%%%%%%%%%%%
Per ogni periodo vengono riportati di solito le seguenti informazioni:
\begin{itemize}
    \item Data inizio,data fine prevista,data di fine attuale ed eventuale giorni di ritardo;
    \item Pianificazione delle attività da svolgere al suo interno con eventuali rischi;
    \item Tempo stimato per poter completare le attività previste (preventivo);
    \item Confronto tra il lavoro svolto e quello preventivato, con annessa analisi dei costi;
    \item Rischi riscontrati durante ogni periodo;
    \item Retrospettiva del periodo per capire cosa mantenere, cosa cambiare e cosa migliorare in futuro
\end{itemize}

I periodi vengono a loro volta divisi in 3 macroperiodi che corrispondono alle revisioni di avanzamento del progetto:
\begin{itemize}
    \item RTB (Requirement and Technology Baseline)
    \item PB (Product Baseline)
    \item CA (Customer Acceptance)
\end{itemize}
Alla fine di ogni collettivo si trovano una revisione del calendario rimanente e un aggiornamento della stima dei costi finali, oltre ad un riepilogo del lavoro svolto.
\subsection{RTB}
\subsubsection{Periodo I}
\begin{itemize}
    \item Inizio: 2024-11-14
    \item Fine prevista: 2024-11-25
    \item Fine attuale: 2024-11-XX
    \item Giorni di ritardo: (X)
\end{itemize}
\paragraph{Pianificazione}
\paragraph{Rischi attesi}
\paragraph{Preventivo} \hspace{1cm} 
\\ \hspace{1cm} \\
Ruoli coinvolti: Responsabile,Amministratore,Verificatore,Analista.
\begin{table}[H]
\centering
\resizebox{1.0\textwidth}{!}{%
    \begin{tabular}{|c|c|c|c|c|c|c|}
        \hline
         \textbf{} & \textbf{Responsabile} &  \textbf{Amministratore} &  \textbf{Verificatore} &  \textbf{Analista} & \textbf{Progettista}  &  \textbf{Programmatore} \\
         \hline \colorrow
           \cellcolor{lightgray}{\textbf{Benin}} & 0 & 0 & 0 & 0 & 0 & 0\\
        \hline
           \textbf{Bourosu} & 0 & 0 & 0 & 0 & 0 & 0\\
        \hline \colorrow
           \cellcolor{lightgray}{\textbf{Gerardin}} & 0 & 0 & 0 & 0 & 0 & 0\\
        \hline 
           \textbf{Gusatto} & 0 & 0 & 0 & 0 & 0 & 0\\
        \hline \colorrow
           \cellcolor{lightgray}{\textbf{Martinelli}} & 0 & 0 & 0 & 0 & 0 & 0\\
        \hline 
           \textbf{Piron} & 0 & 0 & 0 & 0 & 0 & 0\\
        \hline \colorrow
           \cellcolor{lightgray}{\textbf{Zocche}} & 0 & 0 & 0 & 0 & 0 & 0\\
        \hline
    \end{tabular}
}
\caption{Preventivo dell'impegno orario per ruolo di ciascun membro durante il periodo 1}
\end{table}
\paragraph{Consuntivo}
\paragraph{Prospetto orario}
\paragraph{Prospetto economico} \hspace{1cm}

\begin{table}[H]
    \centering
    \begin{tabular}{|c|c|c|c|c|}
            \hline
             \textbf{} & \textbf{Ruolo} &  \textbf{Ore} &  \textbf{Costo} &  \textbf{Differenza}  \\
             \hline \colorrow
               \cellcolor{lightgray}{\textbf{Responsabile}} & 0 & 0 & 0 & 0 \\
            \hline
               \textbf{Amministratore} & 0 & 0 & 0 & 0\\
            \hline \colorrow
               \cellcolor{lightgray}{\textbf{Verificatore}} & 0 & 0 & 0 & 0\\
            \hline 
               \textbf{Analista} & 0 & 0 & 0 & 0\\
            \hline \colorrow
               \cellcolor{lightgray}{\textbf{Progettista}} & 0 & 0 & 0 & 0 \\
            \hline 
               \textbf{Programmatore} & 0 & 0 & 0 & 0 \\
            \hline \colorrow
               \cellcolor{lightgray}{\textbf{Totale preventivo}} & 0 & 0 & 0 & 0 \\
            \hline 
               \textbf{Totale consuntivo} & 0 & 0 & 0 & 0 \\
            \hline
        \end{tabular}
    \caption{Aggiornamenti economici del progetto al termine del periodo I, riflettendo le variazioni tra preventivo e ore effettive di lavoro}
\end{table} 

\paragraph{Rischi occorsi,impatto e la loro mitigazione}
\paragraph{Retrospettiva}
\paragraph{Punto di avanzamento raggiunto e prospettive di completamento}
\hspace{1cm} \\ \\ Barra di avanzamento \\ \\
\begin{tikzpicture}
    \draw[gray, fill=gray!20] (0,0) rectangle (14,1);
    \draw[black, fill=black!70] (0,0) rectangle (10.5,1); % 75% completato
    \node at (7,-0.5) {75\%};
\end{tikzpicture}
\subsubsection{Periodo II}
